% =============================================================================
% Chapter: Query Processing & Optimization
% =============================================================================

\chapter{Query Processing \& Optimization}

\begin{formulabox}[Key Concepts - MEMORIZE!]
\textbf{Evaluation Methods:}
\begin{center}
\begin{tabular}{lp{9cm}}
\toprule
\textbf{Method} & \textbf{Description} \\
\midrule
\textbf{Materialized} & Κάθε operator υπολογίζει πλήρως και αποθηκεύει αποτέλεσμα πριν το δώσει στον επόμενο \\
\textbf{Pipelined} & Tuples περνούν στον επόμενο operator αμέσως μόλις παραχθούν (streaming) \\
\bottomrule
\end{tabular}
\end{center}

\textbf{Join Algorithms Complexity:}
\begin{center}
\begin{tabular}{ll}
Nested Loop Join & $O(n \times m)$ \\
Hash Join & $O(n + m)$ \\
Merge Join & $O(n + m)$ (sorted input) \\
\end{tabular}
\end{center}
\end{formulabox}

% =============================================================================
\section{Query Evaluation Pipeline}

\begin{center}
\begin{tikzpicture}[node distance=1.5cm]
    % Parser
    \node[draw=primary, fill=box-blue, rounded corners=5pt, minimum width=3cm, minimum height=1cm] (parser) {Parser};

    % Optimizer
    \node[draw=primary, fill=box-blue, rounded corners=5pt, minimum width=3cm, minimum height=1cm, right=of parser] (opt) {Optimizer};

    % Executor
    \node[draw=primary, fill=box-blue, rounded corners=5pt, minimum width=3cm, minimum height=1cm, right=of opt] (exec) {Executor};

    % Arrows
    \draw[->, thick] (parser) -- node[above, font=\small] {Query Plan} (opt);
    \draw[->, thick] (opt) -- node[above, font=\small] {Optimized Plan} (exec);

    % Labels below
    \node[anchor=north, font=\footnotesize, text=secondary] at (parser.south) {SQL → Relational Algebra};
    \node[anchor=north, font=\footnotesize, text=secondary] at (opt.south) {Cost Estimation};
    \node[anchor=north, font=\footnotesize, text=secondary] at (exec.south) {Execute Plan};
\end{tikzpicture}
\end{center}

% =============================================================================
\section{Materialized vs Pipelined Evaluation}

\subsection{Materialized Evaluation}

\begin{definitionbox}[Materialized Evaluation]
Κάθε operation υπολογίζει \textbf{πλήρως} το αποτέλεσμά του και το αποθηκεύει (temporary table) πριν περάσει στην επόμενη operation.

\textbf{Πλεονεκτήματα:}
\begin{itemize}
    \item Απλή υλοποίηση
    \item Μπορεί να επαναχρησιμοποιηθεί αποτέλεσμα
\end{itemize}

\textbf{Μειονεκτήματα:}
\begin{itemize}
    \item Χρειάζεται μεγάλο disk I/O
    \item Καθυστέρηση μέχρι να τελειώσει κάθε βήμα
\end{itemize}
\end{definitionbox}

\subsection{Pipelined Evaluation}

\begin{definitionbox}[Pipelined Evaluation]
Τα tuples περνούν στον επόμενο operator \textbf{αμέσως} μόλις παραχθούν (on-the-fly).

\textbf{Πλεονεκτήματα:}
\begin{itemize}
    \item Μειωμένο disk I/O
    \item Χαμηλή latency (αρχίζει να επιστρέφει αποτελέσματα νωρίς)
    \item Μικρότερη χρήση μνήμης
\end{itemize}

\textbf{Μειονεκτήματα:}
\begin{itemize}
    \item Πιο περίπλοκη υλοποίηση
    \item Δεν γίνεται πάντα (π.χ. sorting, hash build)
\end{itemize}
\end{definitionbox}

\begin{center}
\begin{tikzpicture}
    % Materialized
    \node[draw=accent, fill=accent!20, rounded corners=8pt, minimum width=6cm, minimum height=4cm] (mat) at (0,0) {};
    \node[anchor=north, font=\bfseries\color{accent}] at (mat.north) {Materialized};

    % Visualization
    \node[draw=primary, minimum width=1.5cm, minimum height=0.8cm] (r1) at (-1.5,0.5) {R};
    \node[draw=secondary, minimum width=1.5cm, minimum height=0.8cm] (temp) at (0,0.5) {Temp};
    \node[draw=success, minimum width=1.5cm, minimum height=0.8cm] (res1) at (1.5,0.5) {Result};

    \draw[->, thick] (r1) -- (temp);
    \draw[->, thick] (temp) -- (res1);

    \node[font=\small] at (0,-1) {Store complete result};
    \node[font=\small] at (0,-1.5) {before next step};

    % Pipelined
    \node[draw=success, fill=success!20, rounded corners=8pt, minimum width=6cm, minimum height=4cm] (pip) at (8,0) {};
    \node[anchor=north, font=\bfseries\color{success}] at (pip.north) {Pipelined};

    % Visualization with streaming
    \node[draw=primary, minimum width=1.5cm, minimum height=0.8cm] (r2) at (6.5,0.5) {R};
    \node[draw=success, minimum width=1.5cm, minimum height=0.8cm] (res2) at (9.5,0.5) {Result};

    \foreach \i in {0,1,2} {
        \draw[->, thick, secondary] ([xshift=0.8cm]r2.east) ++(0.2*\i, 0.1*\i) -- ([xshift=-0.8cm]res2.west);
    }

    \node[font=\small] at (8,-1) {Pass tuples directly};
    \node[font=\small] at (8,-1.5) {(streaming)};
\end{tikzpicture}
\end{center}

% =============================================================================
\section{Join Algorithms}

\subsection{Nested Loop Join}

\begin{algorithmbox}[Nested Loop Join]
\begin{lstlisting}[language=pseudo]
for each tuple r in R:
    for each tuple s in S:
        if r.key == s.key:
            emit (r, s)
\end{lstlisting}

\textbf{Cost:} $O(|R| \times |S|)$ comparisons

\textbf{Best when:}
\begin{itemize}
    \item One table is very small
    \item Index exists on inner table
\end{itemize}
\end{algorithmbox}

\subsection{Hash Join}

\begin{algorithmbox}[Hash Join]
\textbf{Phase 1: Build}
\begin{lstlisting}[language=pseudo]
for each tuple r in R (smaller table):
    insert r into hash_table on r.key
\end{lstlisting}

\textbf{Phase 2: Probe}
\begin{lstlisting}[language=pseudo]
for each tuple s in S:
    for each r in hash_table.lookup(s.key):
        emit (r, s)
\end{lstlisting}

\textbf{Cost:} $O(|R| + |S|)$

\textbf{Best when:}
\begin{itemize}
    \item Equality join condition
    \item Smaller table fits in memory
\end{itemize}
\end{algorithmbox}

\subsection{Merge Join (Sort-Merge Join)}

\begin{algorithmbox}[Merge Join]
\textbf{Requirement:} Both tables sorted on join key

\begin{lstlisting}[language=pseudo]
// Both R, S sorted on join key
r = first(R), s = first(S)
while not end of R and not end of S:
    if r.key < s.key: r = next(R)
    else if r.key > s.key: s = next(S)
    else:  // r.key == s.key
        emit all matching pairs
        advance pointers
\end{lstlisting}

\textbf{Cost:} $O(|R| + |S|)$ (if already sorted) + sorting cost

\textbf{Best when:}
\begin{itemize}
    \item Data already sorted (index)
    \item Result of ORDER BY
\end{itemize}
\end{algorithmbox}

\subsection{Join Algorithm Comparison}

\begin{center}
\begin{tikzpicture}
    \matrix (m) [
        matrix of nodes,
        nodes={
            draw=primary,
            minimum width=3.5cm,
            minimum height=1.2cm,
            anchor=center,
            font=\small
        },
        column 1/.style={nodes={fill=ntua-blue!20, font=\bfseries}},
        row 1/.style={nodes={fill=secondary!30, font=\bfseries}},
        column sep=-\pgflinewidth,
        row sep=-\pgflinewidth
    ] {
        Algorithm & Best Case & Requirements \\
        Nested Loop & Small inner table & Index on inner \\
        Hash Join & Equality join & Memory for hash table \\
        Merge Join & Sorted data & Pre-sorted input \\
    };
\end{tikzpicture}
\end{center}

% =============================================================================
\section{Solved Exam Problems}

\subsection{2024-25 Exam Question 5 Style}

\begin{examplebox}[Pipelined vs Materialized]
\textbf{Question:} Given the query plan below, identify which operations can be pipelined.

\begin{center}
\begin{tikzpicture}[level distance=1.5cm, sibling distance=4cm]
    \node[draw=primary, fill=box-blue, rounded corners=3pt] {$\pi_{name}$}
        child {
            node[draw=primary, fill=box-blue, rounded corners=3pt] {$\bowtie$}
            child {
                node[draw=secondary, fill=secondary!20, rounded corners=3pt] {$\sigma_{dept='CS'}$}
                child {
                    node[draw=gray, rounded corners=3pt] {Student}
                }
            }
            child {
                node[draw=gray, rounded corners=3pt] {Enrolls}
            }
        };
\end{tikzpicture}
\end{center}

\textbf{Answer:}
\begin{itemize}
    \item $\sigma$ (Selection): \textcolor{success}{Pipelined} - φιλτράρει tuples on-the-fly
    \item $\bowtie$ (Join): Εξαρτάται από algorithm
    \begin{itemize}
        \item Hash Join build phase: \textcolor{accent}{Blocking}
        \item Nested Loop: \textcolor{success}{Pipelined}
    \end{itemize}
    \item $\pi$ (Projection): \textcolor{success}{Pipelined} - transforms tuples immediately
\end{itemize}
\end{examplebox}

\subsection{Cost Estimation}

\begin{examplebox}[Disk I/O Cost]
\textbf{Given:}
\begin{itemize}
    \item Table R: 1000 pages, 10000 tuples
    \item Table S: 500 pages, 5000 tuples
    \item Memory: 50 buffer pages
\end{itemize}

\textbf{Question:} Calculate I/O cost for Nested Loop Join.

\textbf{Solution:}
\[
\text{Cost} = |R| + |R| \times |S| = 1000 + 1000 \times 500 = 501000 \text{ I/Os}
\]

\textbf{With Block Nested Loop (using memory):}
\[
\text{Cost} = |R| + \lceil |R|/B \rceil \times |S| = 1000 + \lceil 1000/48 \rceil \times 500 = 1000 + 21 \times 500 = 11500 \text{ I/Os}
\]

(B = buffer pages for R = 50 - 2 = 48)
\end{examplebox}

% =============================================================================
\section{Query Optimization Rules}

\begin{infobox}[Optimization Heuristics]
\begin{enumerate}
    \item \textbf{Push selections down}: Φιλτράρισε νωρίς, μείωσε δεδομένα
    \item \textbf{Push projections down}: Κράτα μόνο needed columns
    \item \textbf{Combine selections}: $\sigma_{c1}(\sigma_{c2}(R)) = \sigma_{c1 \land c2}(R)$
    \item \textbf{Reorder joins}: Βάλε μικρότερο πίνακα ως outer
    \item \textbf{Use indexes}: Avoid full table scans
\end{enumerate}
\end{infobox}

% =============================================================================
\section{Common Mistakes}

\begin{warningbox}[Συχνά Λάθη στις Εξετάσεις]
\begin{enumerate}
    \item \textbf{Hash Join is always best}:
    \begin{itemize}
        \item[$\times$] Hash Join πάντα καλύτερο
        \item[\checkmark] Χρειάζεται equality condition + memory
    \end{itemize}

    \item \textbf{Pipelined = παντού}:
    \begin{itemize}
        \item Sort, Hash build, Aggregation: \textbf{Blocking} operations
        \item Selection, Projection: Pipelined
    \end{itemize}

    \item \textbf{Merge Join χωρίς sorting}:
    \begin{itemize}
        \item Πρέπει data να είναι sorted on join key
        \item Sorting cost μπορεί να είναι significant
    \end{itemize}
\end{enumerate}
\end{warningbox}

% =============================================================================
\section{Quick Reference}

\begin{center}
\begin{tikzpicture}
    \node[draw=ntua-blue, fill=ntua-blue!10, rounded corners=5pt, minimum width=14cm, minimum height=6cm] (main) {};

    \node[anchor=north west, font=\bfseries\Large\color{ntua-blue}] at ([xshift=5pt, yshift=-5pt]main.north west) {Query Processing Summary};

    \node[anchor=north west, align=left, font=\small] at ([xshift=10pt, yshift=-35pt]main.north west) {
        \textbf{Blocking Operations:}\\
        • Sort, Hash table build, Aggregation (needs all input)\\[5pt]
        \textbf{Pipelined Operations:}\\
        • Selection ($\sigma$), Projection ($\pi$), Nested Loop (outer)\\[5pt]
        \textbf{Join Selection:}\\
        • Equality + Memory → Hash Join\\
        • Sorted data → Merge Join\\
        • Small table / Index → Nested Loop
    };

    \node[draw=secondary, fill=box-blue, rounded corners=5pt, minimum width=5cm, minimum height=1.5cm]
        at ([xshift=-4cm, yshift=1.3cm]main.south) {};
    \node[align=center, font=\footnotesize] at ([xshift=-4cm, yshift=1.3cm]main.south) {
        \textbf{Materialized:} Store full result\\
        \textbf{Pipelined:} Stream tuples
    };

    \node[draw=accent, fill=box-red, rounded corners=5pt, minimum width=5cm, minimum height=1.5cm]
        at ([xshift=2cm, yshift=1.3cm]main.south) {};
    \node[align=center, font=\footnotesize] at ([xshift=2cm, yshift=1.3cm]main.south) {
        \textbf{Optimize by:}\\
        Push $\sigma$ down, use indexes
    };
\end{tikzpicture}
\end{center}
